% generated by tools/settings.py, do not edit
\section{Réglages}
\subsection*{Autres}
\begin{description}
\item[Durée par défaut]  \emph{Choix :} 30 secondes, 1 minute, 2 minutes, 3 minutes, 5 minutes, 10 minutes, 15 minutes, 30 minutes. \emph{Par défaut :} 3 minutes. {\footnotesize\texttt{zxtune.core.plugins.default\_duration}}
\item[Silence detection]  \emph{Choix :} Aucun, 3 secondes, 4 secondes, 5 secondes, 10 secondes, 15 secondes. \emph{Par défaut :} Aucun. {\footnotesize\texttt{zxtune.sound.silencelimit}}
\item[Loops limit]  \emph{Choix :} None, 1 loop, 2 loops, 3 loops, 4 loops, 5 loops. \emph{Par défaut :} None. {\footnotesize\texttt{zxtune.sound.looplimit}}
\item[Fade-in]  \emph{Choix :} Aucun, 1 seconde, 2 secondes, 3 secondes, 5 secondes, 10 secondes. \emph{Par défaut :} Aucun. {\footnotesize\texttt{zxtune.sound.fadein}}
\item[Fade-out]  \emph{Choix :} Aucun, 1 seconde, 2 secondes, 3 secondes, 5 secondes, 10 secondes. \emph{Par défaut :} Aucun. {\footnotesize\texttt{zxtune.sound.fadeout}}
\item[Gain] (plage : 50\% à 200\%) \emph{Par défaut :} 100\%. {\footnotesize\texttt{zxtune.sound.gain}}
\end{description}
\subsection*{AY–3–8910/YM2149F}
\begin{description}
\item[Fréquence d'horloge]  \emph{Choix :} ZX Spectrum (1,7734 MHz), Pentagon (1,75 MHz), Taganrog (3,5 MHz), Atari ST (2 MHz), Amstrad CPC (1 MHz). \emph{Par défaut :} Pentagon (1,75 MHz). {\footnotesize\texttt{zxtune.core.aym.clockrate}}
\item[Type]  \emph{Choix :} AY–3–8910, YM2149F. \emph{Par défaut :} AY–3–8910. {\footnotesize\texttt{zxtune.core.aym.type}}
\item[Interpolation]  \emph{Choix :} Aucune, Performance, Qualité. \emph{Par défaut :} Qualité. {\footnotesize\texttt{zxtune.core.aym.interpolation}}
\item[Disposition des canaux]  \emph{Choix :} ABC, ACB, BAC, BCA, CAB, CBA, Mono. \emph{Par défaut :} ABC. {\footnotesize\texttt{zxtune.core.aym.layout}}
\item[Canaux à rapport cyclique]  \emph{Choix :} Aucun, A, B, C, A, B, A, C, B, C, Tous. \emph{Par défaut :} Aucun. {\footnotesize\texttt{zxtune.core.aym.duty\_cycle\_mask}}
\item[Rapport cyclique] Largeur d'impulsion du son pour les canaux choisis, 50 = onde carrée (plage : 1\% à 99\%) \emph{Par défaut :} 50\%. {\footnotesize\texttt{zxtune.core.aym.duty\_cycle}}
\item[3-channel mixer]  {\footnotesize\texttt{zxtune.sound.mixer.3}}
\end{description}
\subsection*{Processeur Z80 (fichiers AY avec code du lecteur)}
\begin{description}
\item[Fréquence d'horloge]  \emph{Choix :} ZX Spectrum 48K, Pentagon (3,5 MHz), ZX Spectrum 128K (3,5469 MHz), Amstrad CPC (4 MHz), Turbo (7 MHz). \emph{Par défaut :} ZX Spectrum 48K, Pentagon (3,5 MHz). {\footnotesize\texttt{zxtune.core.z80.clockrate}}
\item[Durée d'interruption] Durée d'activité de la demande d'interruption à chaque image (plage : 1 T à 100 T) \emph{Par défaut :} 24 T. {\footnotesize\texttt{zxtune.core.z80.int\_ticks}}
\end{description}
\subsection*{YM2203 (TurboFM)}
\begin{description}
\item[Fréquence d'horloge]  \emph{Choix :} 1,75 MHz, TurboFM (3,5 MHz), NTSC (3,58 MHz), 4 MHz. \emph{Par défaut :} TurboFM (3,5 MHz). {\footnotesize\texttt{zxtune.core.fm.clockrate}}
\end{description}
\subsection*{YM2612/YM3438 (OPN2)}
\begin{description}
\item[Cœur d'émulation]  \emph{Choix :} Nuked (précis, CPU élevé), ymfm (MAME), GPGX (rapide). \emph{Par défaut :} Nuked (précis, CPU élevé). {\footnotesize\texttt{zxtune.core.fm.opn2.core}}
\item[Plateforme]  \emph{Choix :} Mega Drive modèle 1 (VA0-VA6, YM2612), Mega Drive modèle 1 VA7 (YM3438), Mega Drive modèle 2 (YM3438), Genesis 3 (YM3438), YM3438 sans filtre, YMF276 sans filtre, Personnalisé. \emph{Par défaut :} Mega Drive modèle 1 (VA0-VA6, YM2612). {\footnotesize\texttt{zxtune.core.fm.opn2.preset}}
\item[Puce]  \emph{Choix :} YM2612, YM3438, YMF276. \emph{Par défaut :} YM2612. {\footnotesize\texttt{zxtune.core.fm.opn2.chip}}
\item[Effet d'escalier du DAC] Distorsion de croisement du YM2612 discret, audible sur les sons faibles (cœurs Nuked et ymfm) \emph{Par défaut :} activé. {\footnotesize\texttt{zxtune.core.fm.opn2.ladder}}
\item[Filtre passe-bas] Coupure du filtre de la carte, 0 le désactive (plage : 0 Hz à 24000 Hz) \emph{Par défaut :} 3390 Hz. {\footnotesize\texttt{zxtune.core.fm.opn2.lowpass}}
\item[Pente du passe-bas]  \emph{Choix :} 6 dB/oct, 12 dB/oct. \emph{Par défaut :} 6 dB/oct. {\footnotesize\texttt{zxtune.core.fm.opn2.lowpass\_order}}
\item[Blocage du continu] Coupure des condensateurs de sortie, 0 garde la composante continue de la puce (plage : 0 Hz à 200 Hz) \emph{Par défaut :} 5 Hz. {\footnotesize\texttt{zxtune.core.fm.opn2.highpass}}
\item[Largeur stéréo] (plage : 0\% à 100\%) \emph{Par défaut :} 100\%. {\footnotesize\texttt{zxtune.core.fm.opn2.stereo}}
\item[Réponse de sortie]  {\footnotesize\texttt{ui.fm.opn2.response}}
\end{description}
\subsection*{YM2151 (OPM)}
\begin{description}
\item[Cœur d'émulation]  \emph{Choix :} Nuked (précis, CPU très élevé), ymfm (MAME), MAME (ancien, rapide). \emph{Par défaut :} ymfm (MAME). {\footnotesize\texttt{zxtune.core.fm.opm.core}}
\item[Plateforme]  \emph{Choix :} YM3012 (DAC stéréo : X68000, arcade), YM3014 (DAC mono), Personnalisé. \emph{Par défaut :} YM3012 (DAC stéréo : X68000, arcade). {\footnotesize\texttt{zxtune.core.fm.opm.preset}}
\item[Filtre passe-bas] Coupure du filtre de la carte, 0 le désactive (plage : 0 Hz à 24000 Hz) \emph{Par défaut :} 0 Hz. {\footnotesize\texttt{zxtune.core.fm.opm.lowpass}}
\item[Pente du passe-bas]  \emph{Choix :} 6 dB/oct, 12 dB/oct. \emph{Par défaut :} 12 dB/oct. {\footnotesize\texttt{zxtune.core.fm.opm.lowpass\_order}}
\item[Blocage du continu] Coupure des condensateurs de sortie, 0 garde la composante continue de la puce (plage : 0 Hz à 200 Hz) \emph{Par défaut :} 0 Hz. {\footnotesize\texttt{zxtune.core.fm.opm.highpass}}
\item[Largeur stéréo] (plage : 0\% à 100\%) \emph{Par défaut :} 100\%. {\footnotesize\texttt{zxtune.core.fm.opm.stereo}}
\item[Réponse de sortie]  {\footnotesize\texttt{ui.fm.opm.response}}
\end{description}
\subsection*{YM2203/YM2608/YM2610 (OPN)}
\begin{description}
\item[Cœur d'émulation]  \emph{Choix :} ymfm (MAME), MAME (ancien). \emph{Par défaut :} ymfm (MAME). {\footnotesize\texttt{zxtune.core.fm.opn.core}}
\item[Plateforme]  \emph{Choix :} PC-88/PC-98, Personnalisé. \emph{Par défaut :} PC-88/PC-98. {\footnotesize\texttt{zxtune.core.fm.opn.preset}}
\item[Filtre passe-bas] Coupure du filtre de la carte, 0 le désactive (plage : 0 Hz à 24000 Hz) \emph{Par défaut :} 0 Hz. {\footnotesize\texttt{zxtune.core.fm.opn.lowpass}}
\item[Pente du passe-bas]  \emph{Choix :} 6 dB/oct, 12 dB/oct. \emph{Par défaut :} 12 dB/oct. {\footnotesize\texttt{zxtune.core.fm.opn.lowpass\_order}}
\item[Blocage du continu] Coupure des condensateurs de sortie, 0 garde la composante continue de la puce (plage : 0 Hz à 200 Hz) \emph{Par défaut :} 0 Hz. {\footnotesize\texttt{zxtune.core.fm.opn.highpass}}
\item[Largeur stéréo] (plage : 0\% à 100\%) \emph{Par défaut :} 100\%. {\footnotesize\texttt{zxtune.core.fm.opn.stereo}}
\item[Réponse de sortie]  {\footnotesize\texttt{ui.fm.opn.response}}
\end{description}
\subsection*{YM3812/YMF262 (OPL)}
\begin{description}
\item[Cœur d'émulation]  \emph{Choix :} Nuked (précis), AdLibEmu (rapide), MAME (rapide). \emph{Par défaut :} Nuked (précis). {\footnotesize\texttt{zxtune.core.fm.opl.core}}
\item[Plateforme]  \emph{Choix :} AdLib (mono), Sound Blaster 1.x/2.0, Sound Blaster Pro 1/2, Sound Blaster 16, Personnalisé. \emph{Par défaut :} Sound Blaster 16. {\footnotesize\texttt{zxtune.core.fm.opl.preset}}
\item[Filtre passe-bas] Coupure du filtre de la carte, 0 le désactive (plage : 0 Hz à 24000 Hz) \emph{Par défaut :} 0 Hz. {\footnotesize\texttt{zxtune.core.fm.opl.lowpass}}
\item[Pente du passe-bas]  \emph{Choix :} 6 dB/oct, 12 dB/oct. \emph{Par défaut :} 6 dB/oct. {\footnotesize\texttt{zxtune.core.fm.opl.lowpass\_order}}
\item[Blocage du continu] Coupure des condensateurs de sortie, 0 garde la composante continue de la puce (plage : 0 Hz à 200 Hz) \emph{Par défaut :} 0 Hz. {\footnotesize\texttt{zxtune.core.fm.opl.highpass}}
\item[Largeur stéréo] (plage : 0\% à 100\%) \emph{Par défaut :} 100\%. {\footnotesize\texttt{zxtune.core.fm.opl.stereo}}
\item[Réponse de sortie]  {\footnotesize\texttt{ui.fm.opl.response}}
\end{description}
\subsection*{YM2413 (OPLL)}
\begin{description}
\item[Cœur d'émulation]  \emph{Choix :} Nuked (précis), EMU2413 (rapide), MAME (rapide). \emph{Par défaut :} Nuked (précis). {\footnotesize\texttt{zxtune.core.fm.opll.core}}
\item[Filtre passe-bas] Coupure du filtre de la carte, 0 le désactive (plage : 0 Hz à 24000 Hz) \emph{Par défaut :} 0 Hz. {\footnotesize\texttt{zxtune.core.fm.opll.lowpass}}
\item[Pente du passe-bas]  \emph{Choix :} 6 dB/oct, 12 dB/oct. \emph{Par défaut :} 12 dB/oct. {\footnotesize\texttt{zxtune.core.fm.opll.lowpass\_order}}
\item[Blocage du continu] Coupure des condensateurs de sortie, 0 garde la composante continue de la puce (plage : 0 Hz à 200 Hz) \emph{Par défaut :} 0 Hz. {\footnotesize\texttt{zxtune.core.fm.opll.highpass}}
\item[Réponse de sortie]  {\footnotesize\texttt{ui.fm.opll.response}}
\item[Interpolation]  \emph{Par défaut :} désactivé. {\footnotesize\texttt{zxtune.core.dac.interpolation}}
\item[Fréquence d'horloge]  \emph{Choix :} SAM Coupé (8 MHz), Creative Music System (7,16 MHz). \emph{Par défaut :} SAM Coupé (8 MHz). {\footnotesize\texttt{zxtune.core.saa.clockrate}}
\item[Interpolation]  \emph{Choix :} Aucune, Performance, Qualité. \emph{Par défaut :} Performance. {\footnotesize\texttt{zxtune.core.saa.interpolation}}
\end{description}
\subsection*{Écran principal}
\begin{description}
\item[Afficher les informations du morceau] Titre, puces, format du flux, émulation des puces, pilote d'origine et moteur au-dessus des commandes. Aussi basculé par le bouton à gauche des commandes \emph{Par défaut :} désactivé. {\footnotesize\texttt{ui.brief.expanded}}
\item[Afficher l'auteur] Auteur affiché sous le titre du morceau \emph{Par défaut :} désactivé. {\footnotesize\texttt{ui.brief.show\_author}}
\end{description}
\subsection*{Oscilloscope}
\begin{description}
\item[Oscilloscope] Affichage, déclenchement et performances {\footnotesize\texttt{ui.oscilloscope}}
\end{description}
\subsection*{Roland MT-32 (MIDI)}
\begin{description}
\item[Dossier des ROM] Choisir le dossier des dumps de ROM de contrôle et PCM du MT-32 ou CM-32L (entiers ou moitiés de puces, détectés par leur contenu) {\footnotesize\texttt{zxtune.core.mt32.roms\_path}}
\item[Modèle préféré]  \emph{Choix :} Automatique (MT-32 le plus récent d'abord), MT-32, CM-32L / LAPC-I. \emph{Par défaut :} Automatique (MT-32 le plus récent d'abord). {\footnotesize\texttt{zxtune.core.mt32.model}}
\item[Voix de l'oscilloscope]  \emph{Choix :} 9 parties (8 mélodiques + rythme), 32 partiels (synthèse LA). \emph{Par défaut :} 9 parties (8 mélodiques + rythme). {\footnotesize\texttt{zxtune.core.mt32.voices}}
\end{description}
\subsection*{SNES (SPC, SNSF)}
\begin{description}
\item[Emulation engine] ares: cycle accurate, plays stereo as the console. snes\_spc: lighter, without phase inverted surround. Applied from the next track \emph{Choix :} snes\_spc, ares. \emph{Par défaut :} ares. {\footnotesize\texttt{zxtune.core.spc.engine}}
\end{description}
\subsection*{MOS6581/SID}
\begin{description}
\item[Moteur d'émulation]  \emph{Choix :} reSID, reSIDfp. \emph{Par défaut :} reSIDfp. {\footnotesize\texttt{zxtune.core.sid.engine}}
\item[Modèle de puce par défaut] Utilisé si le morceau ne le précise pas \emph{Choix :} MOS6581, MOS8580. \emph{Par défaut :} MOS6581. {\footnotesize\texttt{zxtune.core.sid.model}}
\item[Forcer le modèle de puce] Ignorer la valeur indiquée dans le morceau \emph{Par défaut :} désactivé. {\footnotesize\texttt{zxtune.core.sid.model\_force}}
\item[Modèle de C64 par défaut] \%s (utilisé si le morceau ne le précise pas) \emph{Choix :} PAL, NTSC, NTSC (old), Drean (PAL-N), PAL-M. \emph{Par défaut :} PAL. {\footnotesize\texttt{zxtune.core.sid.clock}}
\item[Forcer le modèle de C64] Ignorer la valeur indiquée dans le morceau \emph{Par défaut :} désactivé. {\footnotesize\texttt{zxtune.core.sid.clock\_force}}
\item[Interpolation]  \emph{Choix :} Aucune, Performance, Qualité. \emph{Par défaut :} Aucune. {\footnotesize\texttt{zxtune.core.sid.interpolation}}
\item[zxtune.core.sid.filter]  \emph{Par défaut :} désactivé. {\footnotesize\texttt{zxtune.core.sid.filter}}
\item[Préréglage du filtre MOS6581]  \emph{Choix :} Custom, JSIDPlay2 reSID: Lightest, JSIDPlay2 reSID: Lighter, JSIDPlay2 reSID: Light, JSIDPlay2 reSID: Average, JSIDPlay2 reSID: Dark, JSIDPlay2 reSID: Darker, JSIDPlay2 reSID: Darkest, VICE default (reSID bias +500 mV), ReSID 6581 (resid-fp fit), Alankila 6581R4AR 3789, Alankila 6581R3 3984 \#1, Alankila 6581R3 3984 \#2, Lord Nightmare 6581R3 4285, Lord Nightmare 6581R3 4485, Lord Nightmare 6581R4 1986S, ZrX 6581R3 0384, ZrX 6581R3 1984, ZrX 6581R3 3684, ZrX 6581R3 3985, ZrX 6581R4AR 2286, Trurl 6581R3 0784, Trurl 6581R3 0486S, Trurl 6581R3 3384, Trurl 6581R3 4885, Trurl 6581R4AR 3789, Trurl 6581R4AR 4486, Nata 6581R3 2083, Grue 6581R4AR 3488, KruLLo, Enigma 6581R3 4885, Enigma 6581R3 1585. \emph{Par défaut :} Custom. {\footnotesize\texttt{zxtune.core.sid.filter\_6581\_preset}}
\item[Préréglage du filtre MOS8580]  \emph{Choix :} Custom, JSIDPlay2 reSID: Light, JSIDPlay2 reSID: Average, JSIDPlay2 reSID: Dark, Trurl 8580R5 1489, Trurl 8580R5 3691. \emph{Par défaut :} Custom. {\footnotesize\texttt{zxtune.core.sid.filter\_8580\_preset}}
\item[Courbe de coupure du filtre]  {\footnotesize\texttt{ui.sid.filter\_curve}}
\item[Digi boost du MOS8580] MOS8580 seulement : rend audibles les échantillons joués par le registre de volume \emph{Par défaut :} désactivé. {\footnotesize\texttt{zxtune.core.sid.digiboost}}
\item[reSID : biais du filtre MOS6581] Moteur reSID, filtre MOS6581 personnalisé seulement : décale la correspondance de la fréquence de coupure (plage : -1000 mV à 1000 mV) \emph{Par défaut :} 0 mV. {\footnotesize\texttt{zxtune.core.sid.filter\_bias}}
\item[reSIDfp : courbe du filtre MOS6581] Moteur reSIDfp, filtre personnalisé seulement : 0 = clair, 100 = sombre \emph{Par défaut :} 50. {\footnotesize\texttt{zxtune.core.sid.filter\_6581\_curve}}
\item[reSIDfp : plage du filtre MOS6581] Moteur reSIDfp, filtre personnalisé seulement : 0 = sombre, 100 = clair \emph{Par défaut :} 50. {\footnotesize\texttt{zxtune.core.sid.filter\_6581\_range}}
\item[reSIDfp : courbe du filtre MOS8580] Moteur reSIDfp, filtre personnalisé seulement : 0 = clair, 100 = sombre \emph{Par défaut :} 50. {\footnotesize\texttt{zxtune.core.sid.filter\_8580\_curve}}
\item[reSIDfp : formes d'onde combinées]  \emph{Choix :} Moyennes, Faibles, Fortes. \emph{Par défaut :} Moyennes. {\footnotesize\texttt{zxtune.core.sid.combined\_waveforms}}
\item[Impulsion+dent de scie du 6581 (reSIDfp)]  \emph{Choix :} Normale (la plupart des 6581), Forte (puces à impulsion+dent de scie forte). \emph{Par défaut :} Normale (la plupart des 6581). {\footnotesize\texttt{zxtune.core.sid.strong\_pulse\_saw}}
\item[Sortie 2SID/3SID]  \emph{Choix :} Stéréo (puces réparties), Mono. \emph{Par défaut :} Stéréo (puces réparties). {\footnotesize\texttt{zxtune.core.sid.multi\_sid\_mono}}
\item[Voie des samples (\$D418)] Oscilloscope : une voie en plus avec les samples du volume principal (digis 4 bits), les trois voies montrent les oscillateurs seuls. Appliqué au morceau suivant \emph{Par défaut :} désactivé. {\footnotesize\texttt{zxtune.core.sid.digi\_lane}}
\end{description}
\subsection*{Android Auto}
\begin{description}
\item[Android Auto] Options de l'écran de la voiture {\footnotesize\texttt{ui.car}}
\end{description}
\subsection*{Sauvegarde de la playlist}
\begin{description}
\item[Sauvegarde de la playlist] Exporter ou importer la playlist, exporter un journal de diagnostic {\footnotesize\texttt{rescue}}
\end{description}
\section{Oscilloscope}
\subsection*{Affichage}
\begin{description}
\item[Fenêtre des formes d'onde] Durée affichée de la forme d'onde de chaque voix (plage : 5 ms à 100 ms) \emph{Par défaut :} 40 ms. {\footnotesize\texttt{ui.oscilloscope.window\_ms}}
\item[Amplification visuelle]  \emph{Choix :} Automatique (par voix), $\times$1, $\times$1,5, $\times$2, $\times$3, $\times$4, $\times$6, $\times$8. \emph{Par défaut :} Automatique (par voix). {\footnotesize\texttt{ui.oscilloscope.gain}}
\item[Périodes affichées au plus (zoom sur les notes aiguës)]  \emph{Choix :} Désactivé (toute la fenêtre), 4 périodes, 8 périodes, 16 périodes, 32 périodes. \emph{Par défaut :} 8 périodes. {\footnotesize\texttt{zxtune.sound.scope.max\_periods}}
\item[Fenêtre des jauges SID] Durée affichée des jauges des oscillateurs et du volume général (plage : 2 ms à 40 ms) \emph{Par défaut :} 10 ms. {\footnotesize\texttt{ui.oscilloscope.gauge\_window\_ms}}
\item[Hauteur de l'oscilloscope] Par rapport au 16:9 de la largeur, le tableau de bord en dessous prend le reste de l'écran (plage : 60\% à 160\%) \emph{Par défaut :} 100\%. {\footnotesize\texttt{ui.oscilloscope.height}}
\item[Hauteur de la réponse du filtre] Part de l'oscilloscope prise par la réponse du filtre SID au-dessus des formes d'onde des voix (plage : 20\% à 80\%) \emph{Par défaut :} 55\%. {\footnotesize\texttt{ui.oscilloscope.filter\_share}}
\item[Voies de batterie AdLib] OPL/OPL2 : voies des batteries du mode rythme (peu de musiques l'utilisent), sinon 9 voix en 3$\times$3. Appliqué au morceau suivant \emph{Par défaut :} désactivé. {\footnotesize\texttt{ui.oscilloscope.opl\_rhythm}}
\item[Formes d'onde sous la réponse du filtre] La vue du filtre SID montre la courbe au-dessus des formes d'onde des voix, sinon la courbe seule \emph{Par défaut :} activé. {\footnotesize\texttt{ui.oscilloscope.filter\_waves}}
\item[Afficher les statistiques techniques] Lignes de l'oscilloscope, de la sortie son, du rendu et du GPU sous le tableau de bord \emph{Par défaut :} activé. {\footnotesize\texttt{ui.oscilloscope.show\_stats}}
\item[Limite d'images par seconde]  \emph{Choix :} Fréquence de l'écran, 120 i/s, 90 i/s, 60 i/s, 30 i/s. \emph{Par défaut :} Fréquence de l'écran. {\footnotesize\texttt{ui.oscilloscope.max\_fps}}
\item[Moteur graphique]  \emph{Choix :} Automatique (Vulkan si disponible), Vulkan, OpenGL ES. \emph{Par défaut :} Automatique (Vulkan si disponible). {\footnotesize\texttt{ui.oscilloscope.backend}}
\item[Pilote Vulkan personnalisé] Choisir le paquet du pilote (zip avec meta.json ou bibliothèque .so, p. ex. Mesa Turnip pour les GPU Adreno) {\footnotesize\texttt{ui.oscilloscope.vulkan\_driver}}
\item[Utiliser le pilote Vulkan personnalisé] Le pilote du système est utilisé si le pilote personnalisé ne se charge pas \emph{Par défaut :} désactivé. {\footnotesize\texttt{ui.oscilloscope.vulkan\_driver.enabled}}
\end{description}
\subsection*{Déclenchement}
\begin{description}
\item[Préréglage]  \emph{Choix :} Par défaut (modèle corrscope, 40 ms), Notes évolutives (60 ms), Chiptune carré/impulsion, Basses, notes graves, Triangle NES, Économie d'énergie, Personnalisé. \emph{Par défaut :} Par défaut (modèle corrscope, 40 ms). {\footnotesize\texttt{ui.oscilloscope.trigger\_preset}}
\item[Largeur de déclenchement] Taille de l'historique, doit dépasser deux périodes des notes les plus graves (plage : 10 ms à 200 ms) \emph{Par défaut :} 40 ms. {\footnotesize\texttt{zxtune.sound.scope.trigger.width}}
\item[Force des fronts] Préférer les fronts montants raides : les nouvelles notes sont mieux saisies, mais les notes tenues sautent davantage (plage : 0\% à 300\%) \emph{Par défaut :} 100\%. {\footnotesize\texttt{zxtune.sound.scope.trigger.edge\_strength}}
\item[Largeur de pente] (plage : 5 \% de la période à 100 \% de la période) \emph{Par défaut :} 25 \% de la période. {\footnotesize\texttt{zxtune.sound.scope.trigger.slope\_width}}
\item[Force de l'historique] Préférer la ressemblance avec les images précédentes (plage : 0\% à 300\%) \emph{Par défaut :} 100\%. {\footnotesize\texttt{zxtune.sound.scope.trigger.buffer\_strength}}
\item[Réactivité de l'historique] (plage : 0\% à 100\%) \emph{Par défaut :} 50\%. {\footnotesize\texttt{zxtune.sound.scope.trigger.responsiveness}}
\item[Atténuation de l'historique] (plage : 10 \% de la période à 200 \% de la période) \emph{Par défaut :} 50 \% de la période. {\footnotesize\texttt{zxtune.sound.scope.trigger.buffer\_falloff}}
\item[Réinitialiser sous la ressemblance] Oublier l'historique aux nouvelles notes si l'onde ne lui ressemble pas (plage : 0\% à 100\%) \emph{Par défaut :} 30\%. {\footnotesize\texttt{zxtune.sound.scope.trigger.reset\_below}}
\item[Vitesse de suppression du continu] (plage : 0\% à 100\%) \emph{Par défaut :} 0\%. {\footnotesize\texttt{zxtune.sound.scope.trigger.mean\_responsiveness}}
\item[Amplification du signe] Passages par zéro raides pour les ondes sans fronts nets (p. ex. triangle) (plage : 0\% à 300\%) \emph{Par défaut :} 0\%. {\footnotesize\texttt{zxtune.sound.scope.trigger.sign\_strength}}
\item[Plage de recherche] (plage : 10 \% de la largeur à 100 \% de la largeur) \emph{Par défaut :} 50 \% de la largeur. {\footnotesize\texttt{zxtune.sound.scope.trigger.diameter}}
\item[Fronts descendants] P. ex. pour les trompettes échantillonnées \emph{Par défaut :} désactivé. {\footnotesize\texttt{zxtune.sound.scope.trigger.falling}}
\item[Suivi de hauteur] Remettre l'historique à l'échelle quand la hauteur change \emph{Par défaut :} activé. {\footnotesize\texttt{zxtune.sound.scope.trigger.pitch\_tracking}}
\item[Recalage au passage par zéro, échantillons] Glisser jusqu'au passage par zéro montant le plus proche, 0 pour désactiver \emph{Par défaut :} 0. {\footnotesize\texttt{zxtune.sound.scope.trigger.post\_radius}}
\item[Recherches complètes]  \emph{Choix :} À chaque image, 60 par seconde, 30 par seconde. \emph{Par défaut :} À chaque image. {\footnotesize\texttt{zxtune.sound.scope.trigger.search\_rate}}
\item[Sous-échantillonnage du déclenchement]  \emph{Choix :} Automatique, Aucun (pleine fréquence), 2, 4. \emph{Par défaut :} Automatique. {\footnotesize\texttt{zxtune.sound.scope.trigger.subsampling}}
\end{description}
\section{Mixage}
\subsection*{Channel A}
\begin{description}
\item[Left] (plage : 0\% à 100\%) \emph{Par défaut :} 100\%. {\footnotesize\texttt{zxtune.sound.mixer.3.0\_0}}
\item[Right] (plage : 0\% à 100\%) \emph{Par défaut :} 0\%. {\footnotesize\texttt{zxtune.sound.mixer.3.0\_1}}
\end{description}
\subsection*{Channel B}
\begin{description}
\item[Left] (plage : 0\% à 100\%) \emph{Par défaut :} 60\%. {\footnotesize\texttt{zxtune.sound.mixer.3.1\_0}}
\item[Right] (plage : 0\% à 100\%) \emph{Par défaut :} 60\%. {\footnotesize\texttt{zxtune.sound.mixer.3.1\_1}}
\end{description}
\subsection*{Channel C}
\begin{description}
\item[Left] (plage : 0\% à 100\%) \emph{Par défaut :} 0\%. {\footnotesize\texttt{zxtune.sound.mixer.3.2\_0}}
\item[Right] (plage : 0\% à 100\%) \emph{Par défaut :} 100\%. {\footnotesize\texttt{zxtune.sound.mixer.3.2\_1}}
\end{description}
\section{Android Auto}
\subsection*{Android Auto}
\begin{description}
\item[Afficher l'oscilloscope / le tableau de bord dans la voiture] Remplace la pochette sur l'écran de la voiture par une visualisation animée pendant la lecture. Lisez l'avertissement légal avant d'activer. \emph{Par défaut :} désactivé. {\footnotesize\texttt{ui.car.visualization}}
\item[Contenu affiché dans la voiture]  \emph{Choix :} Oscilloscope, Tableau de bord des puces, Les deux. \emph{Par défaut :} Oscilloscope. {\footnotesize\texttt{ui.car.visualization.mode}}
\item[Fréquence d'images dans la voiture]  \emph{Choix :} 1 img/s, 2 img/s, 5 img/s, 10 img/s, 15 img/s. \emph{Par défaut :} 5 img/s. {\footnotesize\texttt{ui.car.visualization.fps}}
\item[Avertissement légal] Règles sur les écrans visibles par le conducteur dans votre pays {\footnotesize\texttt{ui.car.legal}}
\end{description}
